\section{导读：从概念火热回到生产力}

本节先建立整期的阅读方式。王鹤既是北京大学助理教授，也是银河通用创始人兼 CTO；他的叙述同时覆盖学术史、技术路线、创业选择和行业伦理。相比只讨论“机器人是否会很快进入家庭”，这期更重要的问题是：具身智能为什么从计算机视觉的边缘方向变成资本宠儿，为什么语言不是智能的唯一入口，为什么 VLM/VLA 数据比 LLM 数据更稀缺，以及为什么产业最终要用生产力而不是 demo 来证明自己。

这期与 EP109 谢晨访谈形成互补。EP109 从仿真和合成数据基础设施讲“数据怎么造”；EP106 从学术流派、视觉-行动闭环和公司落地逻辑讲“为什么需要这类数据、哪些商业叙事会伤害行业”。两期合起来，能把具身智能从论文概念、模型路线、数据 recipe、硬件约束和真实场景 ROI 串成一条完整链路。

\begin{importantbox}{本期核心命题}
具身智能不是“把大模型塞进机器人”这么简单。它首先要求感知、行动和环境反馈形成闭环；其次要求硬件、数据、模型和场景螺旋上升；最终要用可规模化生产力证明自己，而不是用融资额、遥操作 demo 或空泛通用机器人叙事证明自己。
\end{importantbox}

\begin{knowledgebox}{视觉策略说明}
本视频是固定访谈画面，没有 slides、白板或产品演示。正文只使用封面作为来源识别，正文图像全部为自制概念图，用来解释语言与智能、具身智能学术谱系、Perception-Action Loop、数据 recipe、生产力即产品和资本乱象。
\end{knowledgebox}

\subsection{本章小结}

EP106 的主线是：具身智能从视觉学术边缘走向产业中心，但真正的判据不是故事多热，而是能不能把感知-行动闭环、数据飞轮和商业生产力同时跑起来。

